\chapter{离线认知}

\section{维护授权}

离线认知在宿主提供的维护租约内运行。维护租约明确允许的时段、计算资源、模型服务、网络权限、预算、取消条件和抢占规则。宿主决定何时允许维护；Nous Wave 决定在授权范围内处理哪些认知需求。

维护工作可以包括经历片段整理、日记维护、派生表示重建、认知图式复核、跨领域一致性检查、假设探索和服务投影维护。

\section{输入快照}

每次离线认知运行固定输入快照。快照至少记录主体标识、输入对象及其修订、来源和证据版本、访问控制与生命周期、权威提交水位、工作流版本、生产者签名和资源预算。

离线运行期间发生的新提交不自动进入既有快照。运行结果在提交前重新核验当前权威状态。

\section{假设}

离线探索首先产生假设（Hypothesis）。假设包含候选主张、适用范围、输入修订、支持证据、反例证据、不确定性、目标领域、拟议操作、生产者和运行标识。

假设没有跨领域写入权限。提交前执行来源重新核验、共同来源祖先展开、反例搜索、对象版本核验、访问控制核验和清除状态核验；随后由目标领域所有者接受、拒绝、标记冲突或判定过期。

多个模型给出相同判断不自动构成多个独立证据。独立性由来源谱系和外部观察决定。

\section{停止与抢占}

离线认知受到明确预算约束。预算耗尽、重复循环、缺少新证据、预期信息增益过低、目标领域背压、前台任务抢占、权限变化、清除操作和显式取消均可终止运行。运行记录保存停止原因和已产生的候选结果。
